\section*{Bab \ref{ch:g5:irreversible-changes} --- Perubahan yang Tidak Dapat Dibalik}

\begin{solution}{exo:g5:irreversible-changes:1}
Cokelat yang mencair: dapat dibalik (cokelat itu mengeras lagi ketika
didinginkan). Korek api yang \omterm{def:g4:what-a-fire-needs:burning}{terbakar}: tidak dapat. Garam yang
dilarutkan: dapat dibalik (\omterm{def:g3:separating-mixtures:evaporate}{uapkan} airnya). Roti yang dipanggang: tidak
dapat.
\end{solution}

\begin{solution}{exo:g5:irreversible-changes:2}
Perubahan yang memunculkan zat-zat baru. Di dapur: memasak telur,
memanggang kue (juga: memanggang roti, apel yang dipotong menjadi
cokelat).
\end{solution}

\begin{solution}{exo:g5:irreversible-changes:3}
Dengan \omterm{def:g3:separating-mixtures:evaporate}{menguapkan} airnya: garamnya tertinggal di cawan. Melarutkan
adalah \omterm{def:g5:irreversible-changes:reversible}{perubahan dapat balik}.
\end{solution}

\begin{solution}{exo:g5:irreversible-changes:4}
Bau baru, dan gumpalan (padatan) yang muncul di dalam susu.
\end{solution}

\begin{solution}{exo:g5:irreversible-changes:5}
Perubahan di sebelah kiri: gula yang \omterm{def:g2:mixing-and-dissolving:dissolve}{larut} dalam air, es yang mencair,
pasir dan kerikil yang \omterm{def:g2:mixing-and-dissolving:mix}{dicampur}. Panah ganda berarti perubahan itu dapat
berlangsung ke dua arah: perubahan itu dapat dibalik.
\end{solution}

\begin{solution}{exo:g5:irreversible-changes:6}
Minyak menjauhkan \omterm{def:g4:air-a-mixture-of-gases:air}{udara} dan air dari besi rantai: minyak memperlambat
perkaratan.
\end{solution}

\begin{solution}{exo:g5:irreversible-changes:7}
Gelembung gas yang muncul di dalam cairan tanpa dipanaskan.
\end{solution}

\begin{solution}{exo:g5:irreversible-changes:8}
Berguna: memasak makanan, memanggang roti, membakar gas untuk pemanas.
Merugikan: perkaratan, makanan yang basi, kayu yang membusuk.
\end{solution}

\begin{solution}{exo:g5:irreversible-changes:9}
Dinginnya memperlambat \omterm{def:g5:irreversible-changes:chemical-change}{perubahan kimia} yang membuat susu menjadi asam,
sehingga susu lebih awet.
\end{solution}

\begin{solution}{exo:g5:irreversible-changes:10}
Tidak. Gelembung itu uap air, yang kembali menjadi air ketika mendingin:
tidak ada yang baru terbentuk, dan perubahan itu dapat dibalik.
Gelembung saja tidak membuktikan bahwa zat baru telah muncul.
\end{solution}

\begin{solution}{exo:g5:irreversible-changes:11}
Lilin yang menyala melepaskan panas dan cahaya, dan menghasilkan zat-zat
baru: uap air (gelas dingin di atasnya berembun) dan karbon dioksida.
Lilin yang \omterm{def:g4:what-a-fire-needs:burning}{terbakar} tidak kembali. Sebaliknya, lilin yang mencair
mengeras lagi ketika mendingin.
\end{solution}
